\section{Experimental Setup}

\subsection{Benchmarks and Workloads}

The paper uses three workload families, because the contribution is not a
single benchmark trick and the protocols are not interchangeable.

The routing mechanism is evaluated on sparse after-ingest QA holdouts: TOMATO
\cite{tomato} and a frozen motion-heavy MVBench slice \cite{mvbench}. These
are the cleanest settings for testing whether training-free temporal reuse
preserves answers at lower effective fresh-frame budgets.

The first-pass speedup story uses Gemma 4-E4B-4bit on Video-MME
\cite{videomme}, MVBench, and TOMATO at 8-frame settings, plus dev-expansion
cells and deeper frame-count probes. The after-ingest persistent-KV story uses
repeated follow-up questions on short Video-MME clips, where the same video is
queried multiple times and prompt reuse is therefore meaningful.

The streaming and UI studies are kept separate from those benchmark claims.
They are native-rate or deployment-style evidence, not replacements for the
main evaluation.
By design, sparse after-ingest QA uses pixel-diff on sampled decoded frames,
while the companion streaming lane uses H.264 metadata at native frame rate.
The signals are protocol-specific, not interchangeable.

\subsection{Models, Hardware, and Preprocessing}

The Qwen routing and persistent-KV experiments use Qwen2.5-VL-7B-Instruct-4bit
through MLX-VLM on Apple Silicon. The dense wall-clock baseline and the
persistent-KV scaling curve use the same local stack on an M3 Air with 16\,GB
unified memory. The first-pass vision-pruning results use Gemma 4-E4B-4bit on
the same general MLX environment. The deployment-scale streaming evaluation
uses a separate RTSP and screen-recording pipeline with Gemma 4 26B.

For benchmark comparability, the main-path runs use uniform global sampling at
8 frames and square-pad plus resize to 560\,$\times$\,560. The routing planner
operates on decoded RGB frames. Padding-aware reuse accounting is reported
separately from raw padded reuse in the artifact summaries; paper-facing
effective-budget numbers use active-region reuse rather than trivially static
padded borders.

\subsection{Metrics and Reporting Discipline}

We report different metrics for different regimes because the paper would be
misleading otherwise.

For Qwen routing, the paper reports accuracy, dense-versus-cached agreement,
parse failures, and effective fresh frames. For Gemma first-pass pruning, it
reports end-to-end speedup, vision-time reduction, accuracy delta, and the
share$\times$reduction ceiling prediction. For persistent-KV, it reports
follow-up latency, speedup relative to the first query, prefix coverage, and
accuracy delta between baseline and session modes.

All automated main-text assets are generated from checked-in artifact summaries.
If a result is advisory, imported, or still provisional, the manuscript says so
explicitly instead of letting the build silently promote it.

\begin{sectiontodo}{Add the final hardware/software table and release sentence}
This section still needs one compact table with model names, hardware, frame
counts, question protocol, and scorer type, plus a final public-release
sentence once the artifact bundle is frozen.

\textcolor{todoBorder}{Source pointers: \path{docs/benchmark-taxonomy.md},
\path{docs/benchmark-setup.md},
\path{docs/methodology/performance.md},
\path{docs/methodology/preprocessing.md},
\path{research/benchmark_manifests/}.}
\end{sectiontodo}
